% SPDX-License-Identifier: Apache-2.0
% covers: ScanFetch
\datasheet{ScanFetch}{The bus-clock end of a scanout from memory: a line
request taken and \code{LineFetch} started on it.}

\dsfacts{\code{lib/parts/src/scanout.rs}}%
{\code{txhdl\_parts::scanout::ScanFetch<A, WC, LEN>}}%
{\code{//docs:examples}}

\subsection*{Function}

\code{LinePair}, on the pixel clock, asks for each line a line ahead;
the request crosses to the bus clock through \code{ChanCdc}, and
\code{ScanFetch} is what takes it there. It hands the line's address
and length to \code{LineFetch}'s \code{base}, \code{words} and
\code{go}, and takes the next request only once that fetch has started
and finished (issue 151).

\subsection*{Parameters}

\begin{center}\footnotesize
\begin{tabular}{@{}lp{0.68\textwidth}@{}}
\toprule
Parameter & Meaning \\
\midrule
\code{A} & The width of an address, in bits. \\
\code{WC} & The width of \code{LineFetch}'s word count. \\
\code{LEN} & The words in a line, the count every fetch asks for. \\
\bottomrule
\end{tabular}
\end{center}

The tables below use \code{ScanFetch<32, 16, 8>}, the fetcher's end of
\code{ex\_scanpair}.

\dstables{ScanFetch}

\subsection*{Behaviour}

\begin{itemize}
\item \code{at}, \code{count} and \code{start} go to \code{LineFetch}'s
\code{base}, \code{words} and \code{go}: the line's address, \code{LEN},
and a pulse of one cycle.
\item A request is taken from \code{lines} only when no fetch is
running, none is about to start, and none has been started that has
not yet said it is running.
\item The request is cut to \code{A} bits.
\end{itemize}

\subsection*{Verification}

\begin{itemize}
\item \code{ex\_scanpair}, as \code{LinePair}'s sheet says, with this
on the bus clock between the crossing and \code{LineFetch}.
\item \code{//docs:scanpair\_sim\_scanfetch\_scanfetch\_tb\_test} and
\code{//docs:scanpair\_vsim\_scanfetch\_test} simulate the netlist
against that run under nvc and Verilator.
\end{itemize}

\subsection*{Used by}

The example \code{ex\_scanpair}. Nothing on a board holds it yet: the
join with the video peripheral is the rest of issue 151.

\subsection*{Limits}

One line at a time: the next request waits for the fetch before it to
finish, so the pixel side must ask no faster than a line's fetch takes.
It runs on one clock, \code{DefaultClock}.
